% ============================================================================
% CHAPITRE 1 — CONTEXTE, \'ETAT DE L'ART ET PROBL\'EMATIQUE
% Section 1.3 : \'Etat de l'art
% source : partie1/etat_art.tex
% ============================================================================

\section{\'ETAT DE L'ART}
\label{sec:etat_art}
\begin{sloppypar}
La cryptographie institutionnelle ne peut \^etre formalis\'ee qu'en r\'eunissant six corpus th\'eoriques d'ordinaire s\'epar\'es, \`a savoir la th\'eorie de l'information classique~\cite{Shannon48,Shannon49,CoverThomas06}, la th\'eorie de l'information
quantique~\cite{NielsenChuang2010,Wilde2017}, les primitives cryptographiques d'obscurcissement~\cite{Pedersen91,Groth16}, les
d\'efinitions pr\'eliminaires du cadre CRO, les trilemmes d'impossibilit\'e~\cite{Arrow50,Brewer00,FLP85} ainsi que la v\'erification formelle~\cite{moura2015lean,Lean4,Coq,Isabelle}. Nous les parcourons ici dans le but de \emph{pr\'eparer des blocs directement utilisables}. Chaque section ne retient donc que les r\'esultats dont les sections~\ref{sec:clo}, \ref{sec:cro_classique},
\ref{sec:cro_pq}, \ref{sec:sotp_construction} et~\ref{sec:lean4}
auront besoin, en indiquant le r\^ole exact qu'ils jouent dans la d\'emonstration.

L'ordre suivi reprend la logique annonc\'ee au paragraphe~\ref{par:progression_FBESCX}, qui structure l'ensemble de la th\`ese. Les
sous-sections~\ref{subsec:fondements_info}, \ref{subsec:fondements_info_quantique}
et~\ref{subsec:primitives_crypto} rassemblent les outils n\'ecessaires pour
d\'efinir les m\'etriques (\textbf{F}). La sous-section~\ref{subsec:trilemmes}
examine ensuite les paradigmes d'impossibilit\'e conditionnelle dont s'inspireront
l'\'etablissement des bornes (\textbf{B}) puis leur extension
post-quantique (\textbf{E}). La sous-section~\ref{subsec:definitions_preliminaires_cro} fixe le
vocabulaire formel dans lequel seront \'enonc\'ees les constructions con\c{c}ues
pour approcher les bornes (\textbf{S}), avant que la sous-section~\ref{subsec:verification_formelle}
ne pr\'esente l'infrastructure de certification (\textbf{C}). Les fondements
cognitifs de la sous-section~\ref{subsec:fondements_cognitifs} motivent enfin
l'hypoth\`ese de travail sur la capacit\'e $\Delta$, que le programme
exp\'erimental devra mesurer (\textbf{X}).

Deux pr\'ecisions de m\'ethode doivent \^etre apport\'ees d\`es \`a pr\'esent. En premier lieu, les primitives de la
cryptographie institutionnelle sont celles de la cryptographie standard~\cite{Goldreich01,DiffieHellman76}.
Ce qui la distingue des sous-disciplines existantes tient \`a son \emph{objectif}, qui consiste \`a quantifier et \`a optimiser l'utilisabilit\'e
institutionnelle des garanties cryptographiques, en traitant notamment
l'\emph{opposabilit\'e} comme une propri\'et\'e de premi\`ere classe. En second lieu,
les trois m\'etriques ($\Conf$~\ref{def:confidentialite_clo}, $\Rel$~\ref{def:fiabilite_clo}, $\Opp$~\ref{def:opposabilite_clo}) seront d\'efinies ici dans
leur formulation issue de la th\'eorie de l'information, qui autorise le raisonnement
asymptotique. Leurs \'equivalents computationnels, \'etablis au moyen de jeux de
s\'ecurit\'e~\cite{GoldwasserMicali84,GMR88}, sont trait\'es \`a la section~\ref{sec:cro_classique}.
\end{sloppypar}

\begin{boiteArbitre}
L'arbitre ne dispose pas encore d'un langage pour décrire ce qu'il vit. Ce chapitre
lui en fournit les éléments, avec l'entropie pour mesurer ce qu'une preuve révèle,
l'inégalité de Fano pour borner ce qu'un interprète peut
récupérer d'un signal bruité, ou encore les trilemmes CAP et FLP qui lui montrent que son incapacité s'apparente à
d'autres impossibilités déjà connues. Il ignore encore que son cas admet une borne quantitative,
ce que le chapitre suivant lui apprendra.
\end{boiteArbitre}

% ----------------------------------------------------------------------------
\subsection{Fondements en th\'eorie de l'information classique}
\label{subsec:fondements_info}
% ----------------------------------------------------------------------------
\begin{sloppypar}
Nous avons retenu l'entropie, l'information
mutuelle et les in\'egalit\'es classiques de Shannon~\cite{Shannon48,Shannon49} afin
de quantifier confidentialit\'e et opposabilit\'e dans un m\^eme formalisme.
Les objets d\'efinis ci-apr\`es interviennent dans chacune des preuves des chapitres~\ref{chap:cadre_formel} et~\ref{chap:constructions}.

\subsubsection{Entropie et information mutuelle}
\label{subsubsec:entropie_info_mutuelle}

L'entropie de Shannon d'une variable al\'eatoire discr\`ete $X$ sur alphabet
$\mathcal{X}$ est d\'efinie par~\cite{Shannon48,CoverThomas06,Fano61}
\begin{equation}
H(X) = -\sum_{x \in \mathcal{X}} P(x)\log_2 P(x).
\label{eq:entropie_shannon}
\end{equation}
Elle mesure l'incertitude r\'esiduelle sur $X$ avant toute observation.
L'information mutuelle conditionnelle
\begin{equation}
I(X;Y|Z) = H(X|Z) - H(X|Y,Z)
\label{eq:mi_conditionnelle}
\end{equation}
d\'ecrit la r\'eduction d'incertitude sur $X$ apport\'ee par $Y$, une fois $Z$
connu. L'opposabilit\'e $\Opp$ sera d\'efinie \`a partir de cette quantit\'e, plus adapt\'ee que l'entropie brute~\cite{Wyner75,WynerZiv76}.

\subsubsection{Distances et divergences}
\label{subsubsec:distances_divergences}

La distance de variation totale entre deux distributions $P$ et $Q$, donn\'ee par
\begin{equation}
\TV(P,Q) = \tfrac{1}{2}\sum_x |P(x) - Q(x)|
\label{eq:tv_distance}
\end{equation}
quantifie la distinguabilit\'e au sens du pire adversaire d\'eterministe~\cite{CoverThomas06}. La
divergence de Kullback-Leibler, d\'efinie par
\begin{equation}
\KL(P \| Q) = \sum_x P(x)\log_2\frac{P(x)}{Q(x)}
\label{eq:kl_divergence}
\end{equation}
majore $\TV$ via l'in\'egalit\'e de Pinsker~\cite{Pinsker1964}, que nous utiliserons pour borner
la confidentialit\'e.

\subsubsection{L'in\'egalit\'e de Pinsker}
\label{subsubsec:pinsker_classique}

\begin{lemma}[In\'egalit\'e de Pinsker~\cite{Pinsker1964,CoverThomas06}]\label{lem:pinsker}
Pour toutes distributions $P$ et $Q$ sur un m\^eme alphabet fini, on a
\begin{equation}
\TV(P,Q)^2 \leq c_0 \cdot \KL(P \| Q),
\qquad c_0 = \frac{1}{2\ln 2} \approx 0{,}721.
\label{eq:pinsker}
\end{equation}
\end{lemma}

\noindent Dans la preuve du Trilemme CRO \`a la section~\ref{sec:cro_classique}, cette in\'egalit\'e relie la
confidentialit\'e $\Conf = 1 - \max_{m_0,m_1} \TV(P_{\Sigma|m_0},
P_{\Sigma|m_1})$ \`a l'information mutuelle $I(M;\Sigma|C)$ via
l'hypoth\`ese de $\beta$-r\'egularit\'e. Elle permet ainsi de convertir un
\'ecart entre distributions en une borne sur le produit $\Conf \cdot \Rel \cdot
\Opp$~\cite{Tsybakov2009}.

\subsubsection{Le lemme de Fano}
\label{subsubsec:fano_classique}

\begin{lemma}[Lemme de Fano~\cite{Fano61,CoverThomas06}]\label{lem:fano}
Soit $X$ une variable al\'eatoire sur un alphabet fini de cardinalit\'e $k$,
$\hat{X}$ un estimateur et $p_e = \Pr[X \neq \hat{X}]$. On a alors
\begin{equation}
H(X|Y) \leq h(p_e) + p_e \log_2(k-1),
\label{eq:fano_inequality}
\end{equation}
o\`u $h(p) = -p\log_2 p - (1-p)\log_2(1-p)$ est la fonction d'entropie
binaire.
\end{lemma}

Une fiabilit\'e \'elev\'ee (faible $p_e$) implique donc une information mutuelle \'elev\'ee.
La contrapos\'ee est tout aussi utile, puisque si $H(X|Y)$ reste grande, aucun
estimateur ne peut retrouver $X$ \`a partir de $Y$ avec pr\'ecision~\cite{HanVerdu}. Dans notre
cadre, il en d\'ecoule directement pour le Trilemme CRO l'in\'egalit\'e
\begin{equation}
1 - \Rel \;\geq\; \frac{H(M|C) - I(M;\Sigma|C) - 1}{\log_2 |\mathcal{M}|}.
\label{eq:rel_fano_bound}
\end{equation}

\subsubsection{L'in\'egalit\'e de traitement des donn\'ees}
\label{subsubsec:dpi}

\begin{lemma}[Data Processing Inequality~\cite{CoverThomas06,Lindblad}]\label{lem:dpi}
Pour toute cha\^ine de Markov $X \to Y \to Z$, on a
\begin{equation}
I(X;Z) \;\leq\; I(X;Y).
\label{eq:dpi}
\end{equation}
\end{lemma}

\noindent Ce r\'esultat, rarement mobilis\'e en
cryptographie, intervient directement dans notre analyse~\cite{Smith09,Alvim12}. La
transformation $M \to (\Sigma, C) \to V = \mathsf{Interpret}(\Sigma, C, K)$ forme
une cha\^ine de Markov d\`es que les ressources propres $K$ de l'institution sont
ind\'ependantes du message sachant la preuve et le contexte, que le verdict soit
d\'eterministe ou non. L'in\'egalit\'e \eqref{eq:dpi} garantit alors que
$I(M;V|C) \leq I(M;\Sigma|C)$, autrement dit ce que le verdict apprend du \emph{message} ne peut exc\'eder
ce que la preuve en contient. Cette in\'egalit\'e fournit l'un des deux termes de la
d\'ecomposition qui fonde, \`a la section~\ref{sec:cro_classique}, l'Hypoth\`ese
d'extraction born\'ee (Proposition~\ref{prop:fondement_extraction})~\cite{ClarksonSchneider10}.

\subsubsection{Information flow quantitatif}
\label{subsubsec:qif}

La litt\'erature sur le \emph{Quantitative Information Flow} (QIF)~\cite{Clarkson05,Smith09,Alvim12,Alvim18}
fournit un cadre formel pour mesurer les fuites d'information dans les syst\`emes
de s\'ecurit\'e. La mesure $g$-leakage~\cite{Alvim12} \'etend l'information mutuelle
en pond\'erant chaque fuite par son effet sur les gains de l'adversaire. Notre d\'efinition
de l'opposabilit\'e, con\c{c}ue comme un ratio de r\'eduction d'incertitude, s'en inspire directement.

\subsubsection{Capacit\'e de canal et codage source-canal}
\label{subsubsec:capacite_canal}

Le th\'eor\`eme de codage source-canal~\cite{Shannon48,CoverThomas06} \'etablit que
la transmission fiable d'une source $X$ sur un canal de capacit\'e $C$ est possible
si et seulement si $H(X) < C$. Transpos\'e au contexte institutionnel, ce r\'esultat conduit \`a
mod\'eliser l'interpr\`ete comme un canal de capacit\'e born\'ee $\Delta$, qui limite
la quantit\'e d'information qu'il peut extraire~\cite{WynerZiv76,Maurer1999}.
\end{sloppypar}

% ----------------------------------------------------------------------------
\subsection{Fondements en th\'eorie de l'information quantique}
\label{subsec:fondements_info_quantique}
% ----------------------------------------------------------------------------
\begin{sloppypar}
L'extension du Trilemme CRO
aux contextes quantiques n\'ecessite des outils de th\'eorie de l'information
quantique qui g\'en\'eralisent les notions classiques~\cite{NielsenChuang2010,Wilde2017,Holevo1973}. Cette section introduit
les concepts fondamentaux mobilis\'es \`a la section~\ref{sec:cro_pq}.

\subsubsection{\'Etats quantiques et entropie de von Neumann}
\label{subsubsec:von_neumann}

Un \'etat quantique est repr\'esent\'e par une matrice de densit\'e $\rho$, c'est-\`a-dire un
op\'erateur positif semi-d\'efini de trace unit\'e sur un espace de Hilbert $\mathcal{H}$~\cite{NielsenChuang2010}.
L'entropie de von Neumann, qui g\'en\'eralise l'entropie de Shannon, est d\'efinie par
\begin{equation}
S(\rho) = -\mathrm{Tr}(\rho \log \rho).
\label{eq:von_neumann}
\end{equation}
Pour un \'etat pur $|\psi\rangle$, l'entropie est nulle ($S(|\psi\rangle\langle\psi|) = 0$). Elle
vaut $S(\rho) = \log_2 d$ sur l'\'etat compl\`etement m\'elang\'e $\rho = I/d$. Lorsque
$\rho$ est diagonale dans la base computationnelle, on retrouve $S(\rho) = H(X)$, o\`u $X$
suit la distribution form\'ee par les probabilit\'es diagonales. La d\'efinition quantique
contient donc le cadre classique comme cas particulier~\cite{Wilde2017,Hiai1991}.

\subsubsection{Information mutuelle quantique}
\label{subsubsec:info_mutuelle_quantique}

Pour un \'etat bipartite $\rho_{AB}$, l'information mutuelle quantique s'\'ecrit
\begin{equation}
I(A;B)_\rho = S(\rho_A) + S(\rho_B) - S(\rho_{AB}).
\label{eq:qmi}
\end{equation}
L'information mutuelle quantique ob\'eit elle aussi \`a une in\'egalit\'e de traitement
des donn\'ees~\cite{Lindblad}, puisque l'application d'un canal CPTP $\Lambda$ au
syst\`eme $B$ ne peut que la r\'eduire,
\begin{equation}
I(A;B)_\rho \geq I(A; \Lambda(B))_\rho.
\label{eq:qdpi}
\end{equation}

\subsubsection{Information de Holevo}
\label{subsubsec:holevo}

L'information de Holevo borne l'information classique accessible depuis
un encodage quantique~\cite{Holevo1973,Schumacher1997}. Pour un ensemble $\{P_X(x), \rho_x\}$ d'\'etat moyen
$\bar{\rho} = \sum_x P_X(x)\rho_x$, elle vaut
\begin{equation}
\chi = S(\bar{\rho}) - \sum_x P_X(x) S(\rho_x).
\label{eq:holevo}
\end{equation}
Le th\'eor\`eme de Holevo \'etablit que l'information accessible $I_{\mathrm{acc}}$
par toute mesure POVM satisfait $I_{\mathrm{acc}} \leq \chi$. L'extension
post-quantique de l'opposabilit\'e s'appuie directement sur cette borne. Quels que
soient les outils quantiques dont il dispose, un interpr\`ete institutionnel n'extrait
en effet du \textit{proof state} quantique qu'au plus $\chi$ bits d'information~\cite{Winter}.

\subsubsection{Distance de trace et fid\'elit\'e}
\label{subsubsec:distance_trace_fidelite}

La distance de trace quantique, qui g\'en\'eralise la variation totale~\cite{NielsenChuang2010}, est d\'efinie par
\begin{equation}
\|\rho - \sigma\|_1 = \mathrm{Tr}\sqrt{(\rho-\sigma)^\dagger(\rho-\sigma)}.
\label{eq:trace_distance}
\end{equation}
Pour des \'etats classiques (diagonaux), $\|\rho - \sigma\|_1 = 2\TV(P,Q)$.
La fid\'elit\'e quantique~\cite{Jozsa}
\begin{equation}
F(\rho, \sigma) = \|\sqrt{\rho}\sqrt{\sigma}\|_1
\label{eq:fidelite}
\end{equation}
est reli\'ee \`a la distance de trace par les in\'egalit\'es de Fuchs-van de Graaf~\cite{FuchsGraaf}
\begin{equation}
1 - \tfrac{1}{2}\|\rho - \sigma\|_1 \leq F(\rho, \sigma) \leq \sqrt{1 - \tfrac{1}{4}\|\rho - \sigma\|_1^2}.
\label{eq:fuchs_vdg}
\end{equation}

\subsubsection{In\'egalit\'e de Pinsker quantique}
\label{subsubsec:pinsker_quantique}

\begin{lemma}[In\'egalit\'e de Pinsker quantique~\cite{Wilde2017,Hiai1991}]\label{lem:pinsker_quantique_etat_art}
Pour tous \'etats quantiques $\rho$ et $\sigma$, on a
\begin{equation}
\|\rho - \sigma\|_1^2 \leq 2 D(\rho\|\sigma),
\label{eq:qpinsker}
\end{equation}
o\`u $D(\rho\|\sigma) = \mathrm{Tr}(\rho\log\rho - \rho\log\sigma)$ est l'entropie
relative quantique.
\end{lemma}

Cette in\'egalit\'e sera utilis\'ee dans la section~\ref{sec:cro_pq} (Lemme~\ref{lem:pinsker_quantique}) pour borner la
confidentialit\'e quantique $\Conf^q$ en fonction de l'information mutuelle
quantique $I(M;\Sigma|C)$~\cite{Ohya1993}.

\subsubsection{Lemme de gentle measurement}
\label{subsubsec:gentle_measurement}

\begin{lemma}[Gentle measurement~\cite{Wilde2017,Winter}]\label{lem:mesure_douce_etat_art}
Soit $\rho$ un \'etat quantique et $\{M_m\}$ une POVM. Si l'issue $m$ se produit
avec probabilit\'e $p_m = \mathrm{Tr}(M_m\rho) \geq 1 - \delta$, alors l'\'etat
post-mesure satisfait
\begin{equation}
\|\rho - \rho'\|_1 \leq 2\sqrt{\delta}.
\label{eq:gentle}
\end{equation}
\end{lemma}

Ce lemme capture l'id\'ee qu'une mesure \guillemotleft~douce~\guillemotright (haute probabilit\'e) perturbe
peu l'\'etat quantique~\cite{OgawaNagaoka}. Dans le contexte institutionnel post-quantique, il borne
la perte d'interpr\'etabilit\'e $\eta_q$ induite par les effets de canal quantique.

\subsubsection{In\'egalit\'e de Fano quantique}
\label{subsubsec:fano_quantique}

\begin{lemma}[Lemme de Fano quantique~\cite{Wilde2017,Winter}]\label{lem:fano_quantique_etat_art}
Soient $X$ encod\'e dans l'\'etat quantique $\rho_X$ et $\hat{X}$ le r\'esultat d'une
mesure POVM avec probabilit\'e d'erreur $p_e$. Alors
\begin{equation}
S(X|\hat{X}) \leq h(p_e) + p_e \log_2(|\mathcal{X}| - 1).
\label{eq:qfano}
\end{equation}
\end{lemma}

Cette g\'en\'eralisation du lemme de Fano classique est utilis\'ee dans la section~\ref{sec:cro_pq} (Lemme~\ref{lem:fano_quantique}) pour borner la fiabilit\'e quantique $\Rel^q$~\cite{Hayashi}.

\subsubsection{Canaux quantiques (CPTP)}
\label{subsubsec:canaux_quantiques}

Un canal quantique est une application compl\`etement positive et pr\'eservant
la trace (CPTP) $\Lambda~: \mathcal{L}(\mathcal{H}_A) \to \mathcal{L}(\mathcal{H}_B)$~\cite{NielsenChuang2010}.
Par le th\'eor\`eme de dilatation de Stinespring~\cite{Stinespring}, tout canal admet une repr\'esentation
unitaire avec environnement auxiliaire de la forme
\begin{equation}
\Lambda(\rho) = \mathrm{Tr}_E\left[U (\rho \otimes |0\rangle\langle 0|_E) U^\dagger\right].
\label{eq:stinespring}
\end{equation}

Trois mod\`eles de canaux sont pertinents pour l'analyse post-quantique~\cite{Wilde2017}~:
\begin{itemize}
\item \textbf{Canal d\'epolarisant}, d\'efini par $\Lambda_{\mathrm{dep}}(\rho) = (1-p)\rho + \frac{p}{d}I_d$
\item \textbf{Canal d'amortissement d'amplitude} $\Lambda_{\mathrm{AD}}(\rho)$, qui mod\'elise la dissipation d'\'energie
\item \textbf{Canal de d\'ephasage}, d\'efini par $\Lambda_{\mathrm{deph}}(\rho) = (1-\lambda)\rho + \lambda Z\rho Z$
\end{itemize}

Ces canaux mod\'elisent les effets de bruit quantique qui d\'egradent
l'interpr\'etabilit\'e des preuves dans les syst\`emes hybrides ou quantiques~\cite{Zurek2003,Wootters1982}.

\subsubsection{Impossibilit\'es quantiques fondamentales}
\label{subsubsec:impossibilites_quantiques}

Deux r\'esultats d'impossibilit\'e quantique \'eclairent les contraintes du Trilemme PQ-CRO.

\textbf{No-cloning theorem.} Wootters et Zurek~\cite{Wootters1982} ont prouv\'e qu'un \'etat quantique inconnu ne peut \^etre copi\'e parfaitement. Cette contrainte interdit la duplication arbitraire des preuves quantiques.

\textbf{Impossibilit\'e du bit commitment quantique.} Mayers~\cite{Mayers1997} et Lo-Chau~\cite{Lo1997} ont montr\'e que, dans le mod\`ele quantique non-relativiste, aucun protocole de bit commitment ne garantit simultan\'ement le secret et la liaison. C'est pourquoi les engagements post-quantiques reposent sur des hypoth\`eses calculatoires~\cite{Unruh2016}.
\end{sloppypar}

\begin{boiteArbitre}
L'arbitre retient l'essentiel. Même équipé d'outils quantiques optimaux,
il ne tirera jamais d'un état de preuve quantique plus de $\chi$ bits.
Cette borne découle des lois de la physique, si bien qu'aucun progrès
des instruments ne la repoussera.
\end{boiteArbitre}

% ----------------------------------------------------------------------------
\subsection{Primitives cryptographiques d'obscurcissement}
\label{subsec:primitives_crypto}
% ----------------------------------------------------------------------------
\begin{sloppypar}
Les constructions SOTP du
chapitre~\ref{chap:constructions} font appel \`a trois primitives. Nous les pr\'esentons
bri\`evement en nous limitant aux propri\'et\'es que ces constructions exploitent,
sans viser l'exhaustivit\'e des sch\'emas.

\subsubsection{Engagements de Pedersen}
\label{subsubsec:pedersen}

Un engagement de Pedersen~\cite{Pedersen91} sur un groupe cyclique $G$ de
g\'en\'erateurs $g, h$ est le couple $\mathsf{Com}(m; r) = g^m h^r$, o\`u $m$
est le message secret et $r$ un al\'ea. Le sch\'ema est parfaitement occultant au sens
de la th\'eorie de l'information, puisque pour tout $m' \neq m$ il existe $r'$ tel que
$\mathsf{Com}(m'; r') = \mathsf{Com}(m; r)$. Sa propri\'et\'e de liaison n'est en revanche
que calculatoire, car produire $m' \neq m$ de m\^eme image revient \`a r\'esoudre le probl\`eme
du logarithme discret~\cite{Schnorr91}. La famille SOTP confie \`a ces engagements un double r\^ole.
Ils encodent le contexte institutionnel $C$ de mani\`ere v\'erifiable, au b\'en\'efice de $\Rel$,
tout en pr\'eservant la confidentialit\'e des donn\'ees sous-jacentes, au b\'en\'efice
de $\Conf$~\cite{ChaumPedersen92}.

\subsubsection{Arguments succincts zero-knowledge (zk-SNARKs)}
\label{subsubsec:zksnarks}

Pour un circuit $\mathcal{C}$, une entr\'ee publique $x$ et un t\'emoin secret $w$,
un zk-SNARK~\cite{Groth16,BCTV14,BenSasson2019} atteste que $\mathcal{C}(x, w)
= 1$ au moyen d'une preuve $\pi$ de taille constante ($\approx 128$--$256$ octets).
Le sch\'ema se d\'efinit par sa compl\'etude, qui garantit qu'un prouveur
honn\^ete convainc toujours le v\'erificateur, par sa solidit\'e (soundness) sous les
hypoth\`eses q-SDH et q-PKE~\cite{Gennaro13,Parno13} ainsi que par la propri\'et\'e de
zero-knowledge parfait, en vertu de laquelle le v\'erificateur
n'apprend rien sur $w$ au-del\`a de la validit\'e de l'assertion~\cite{GMR85}.
SOTP-High recourt \`a ces preuves pour maximiser $\Opp$ dans le r\'egime \`a haute
entropie contextuelle ($H(\mathcal{C}) > \lambda/\log\lambda$)~\cite{PLONK20,Wahby18,Campanelli2019}.

\subsubsection{Syst\`emes de preuves avanc\'es}
\label{subsubsec:systemes_preuves}

Plusieurs syst\`emes de preuves compl\`etent les zk-SNARKs dans la palette disponible.

\textbf{Bulletproofs.} B\"unz et al.~\cite{Bunz19} ont introduit des preuves \`a divulgation nulle sans trusted setup, avec des tailles logarithmiques en la taille du t\'emoin. Leur construction repose sur un argument de produit scalaire (\emph{inner-product argument}) et sur la difficult\'e du logarithme discret~; elle n'est donc pas post-quantique.

\textbf{STARKs.} Ben-Sasson et al.~\cite{BenSasson18} ont d\'evelopp\'e des preuves transparentes (sans trusted setup) et post-quantiques, bas\'ees sur des codes correcteurs d'erreurs et la technique FRI (Fast Reed-Solomon Interactive Oracle Proof of Proximity).

\textbf{Fiat-Shamir.} La transformation de Fiat et Shamir~\cite{FiatShamir87} convertit les preuves interactives en preuves non-interactives via l'application de hachage. Sa s\'ecurit\'e dans le mod\`ele de l'oracle al\'eatoire quantique a \'et\'e analys\'ee, pour les signatures, par Kiltz, Lyubashevsky et Schaffner~\cite{Kiltz2017}.

\subsubsection{La transformation de Johnson-Lindenstrauss}
\label{subsubsec:jl}

La transformation de Johnson-Lindenstrauss~\cite{JohnsonLindenstrauss84,DasguptaGupta03}
garantit qu'une famille de $n$ vecteurs de $\mathbb{R}^d$
peut \^etre projet\'ee dans $\mathbb{R}^k$ avec $k = O(\log n / \epsilon^2)$
en pr\'eservant toutes les distances euclidiennes \`a un facteur $1 \pm \epsilon$
pr\`es~\cite{Vempala05,Indyk98,Indyk06}. Dans SOTP-Base (r\'egime de transition), cette transformation compresse
l'espace contextuel, r\'eduisant la complexit\'e des engagements sans
sacrifier les distances s\'emantiques~\cite{DasguptaKumar14,Larsen17,Gionis99}.

\subsubsection{Fonctions de hachage et filtres de Bloom}
\label{subsubsec:hachage}

Les fonctions de hachage cryptographique~\cite{Damgard87} servent aux engagements et \`a la transformation Fiat-Shamir. Les filtres de Bloom~\cite{Bloom70} fournissent des structures de donn\'ees probabilistes pour la v\'erification d'appartenance avec contr\^ole du taux de faux positifs, pertinents pour les contextes institutionnels \`a grande cardinalit\'e.

\subsubsection{Chiffrement homomorphe}
\label{subsubsec:fhe}

Le chiffrement compl\`etement homomorphe (FHE), introduit par Gentry~\cite{Gentry09} et am\'elior\'e par Brakerski et Vaikuntanathan~\cite{BrakerskiVaikuntanathan11}, permet le calcul sur donn\'ees chiffr\'ees. Les sch\'emas bas\'es sur LWE~\cite{Regev2009} sont consid\'er\'es comme post-quantiques et utilis\'es dans les protocoles de calcul multipartite s\'ecuris\'e~\cite{Yao82,GMW87,Asharov12}.
\end{sloppypar}

% ----------------------------------------------------------------------------
\subsection{Fondements cognitifs~: rationalit\'e born\'ee et param\`etre~$\Delta$}
\label{subsec:fondements_cognitifs}
% ----------------------------------------------------------------------------
\begin{sloppypar}
Le choix de la capacit\'e
d'extraction institutionnelle $\Delta$ et de l'hypoth\`ese de r\'egularit\'e $\beta$
s'appuie sur des r\'esultats \'etablis en sciences cognitives et en th\'eorie
de l'information.

\subsubsection{Rationalit\'e born\'ee et capacit\'e de traitement}
\label{subsubsec:rationalite_bornee}

Selon Herbert Simon~\cite{Simon82,Simon1955}, les d\'ecideurs humains, arbitres compris,
ne recherchent pas l'optimum. Soumis \`a de s\'ev\`eres contraintes cognitives, ils \emph{satisficent},
c'est-\`a-dire qu'ils retiennent une solution jug\'ee suffisamment bonne en raisonnant sur des
mod\`eles mentaux simplifi\'es \`a $O(10)$--$O(100)$ \'el\'ements distinguables. George Miller~\cite{Miller56}
a estim\'e la capacit\'e de la m\'emoire de travail \`a $7 \pm 2$
\emph{chunks} d'information, chiffre que Cowan~\cite{Cowan2001} ram\`ene \`a environ
$4$ \'el\'ements. Pour une t\^ache d'interpr\'etation de preuve cryptographique, un ordre
de grandeur s'obtient en retenant $\approx 3$ bits par \'el\'ement et la borne haute de Miller,
ce qui donne

\begin{equation}
\Delta_{\text{individu}} \;\lesssim\; 7 \times 3 \;=\; 21 \;\text{bits}.
\label{eq:delta_individu}
\end{equation}
Ericsson et Kintsch~\cite{Ericsson95} ont montr\'e que les experts acc\`edent, via
la m\'emoire \`a long terme, \`a $10^3$--$10^4$ sch\'emas, ajoutant
$\log_2(10^4) \approx 13$ bits d'indexation. La d\'elib\'eration institutionnelle
collective amplifie encore cette capacit\'e~\cite{Argyris1978}. Faute de quantification
dans la litt\'erature, nous retenons pour l'institution l'ordre de grandeur
\begin{equation}
\Delta_{\text{institution}} \;\sim\; \text{quelques dizaines de bits}.
\label{eq:delta_institution}
\end{equation}
Ces estimations cumulent des grandeurs h\'et\'erog\`enes, \`a savoir la capacit\'e
de la m\'emoire de travail et la taille d'un r\'epertoire de sch\'emas. Elles ne valent
donc pas d\'erivation et servent seulement \`a motiver l'\emph{hypoth\`ese de travail}
d'une capacit\'e de l'ordre de $20$ \`a $80$ bits, qu'il reviendra au programme
exp\'erimental du chapitre~\ref{chap:experimentation} de mesurer.

\subsubsection{Biais cognitifs et d\'ecision}
\label{subsubsec:biais_cognitifs}

Kahneman~\cite{Kahneman11} a document\'e les biais syst\'ematiques affectant
la d\'ecision humaine, parmi lesquels l'aversion \`a la perte, l'ancrage et la disponibilit\'e. Ces biais
modulent la capacit\'e d'extraction effective, si bien qu'un juge peut disposer de
$\Delta$ bits de capacit\'e mais n'en utiliser qu'une fraction selon le contexte
\'emotionnel ou proc\'edural.

\subsubsection{Extension quantique~: $\Delta_q$ et $\eta_q$}
\label{subsubsec:extension_quantique_delta}

Dans le contexte post-quantique, la capacit\'e d'extraction institutionnelle
doit \^etre adapt\'ee. Nous la d\'efinissons par
\begin{equation}
\Delta_q = \Delta + \delta_{\mathrm{enh}} - \delta_{\mathrm{degr}},
\label{eq:delta_q}
\end{equation}
o\`u $\delta_{\mathrm{enh}}$ repr\'esente l'am\'elioration potentielle due aux
outils quantiques (analyse quantique assist\'ee) et $\delta_{\mathrm{degr}}$
repr\'esente la d\'egradation due \`a la complexit\'e accrue des preuves PQC~\cite{Chatterjee2021}.

Le coefficient de perte d'interpr\'etabilit\'e quantique $\eta_q$ mesure la
d\'egradation s\'emantique induite par un canal quantique $\Lambda$ sur le
support de la preuve. Il s'agit du plus petit r\'eel tel que
\begin{equation}
I(V;\Lambda(\Sigma)|C) \geq I(V;\Sigma|C) - \eta_q\, H(C)
\label{eq:eta_q}
\end{equation}
(D\'efinition~\ref{def:perte_interp_quantique}). Ces param\`etres sont analys\'es en d\'etail \`a la section~\ref{sec:cro_pq}.

\subsubsection{Justification issue de la th\'eorie de l'information pour $\Delta$}
\label{subsubsec:justification_delta}

L'interpr\`ete institutionnel (arbitre) poss\`ede une connaissance pr\'ealable $K$
non captur\'ee dans le contexte formel $C$. Comme il n'acc\`ede au message que par la
preuve, on a $K \perp M \mid (\Sigma, C)$. La r\`egle de cha\^inage appliqu\'ee de deux
mani\`eres \`a $I(V;M,\Sigma|C)$ donne alors
\begin{equation}
I(V;\Sigma|C) \;=\; I(V;M|C) \;+\; I(V;\Sigma|M,C),
\label{eq:chaining_knowledge}
\end{equation}
et l'in\'egalit\'e de traitement des donn\'ees sur $M \to (\Sigma,C) \to V$ donne
$I(V;M|C) \leq I(M;\Sigma|C)$. Si l'institution tire de la preuve au plus
$\Delta$ bits \emph{au-del\`a du message}, $I(V;\Sigma|M,C) \leq \Delta$, d'o\`u
\begin{equation}
I(V;\Sigma|C) \;\leq\; I(M;\Sigma|C) + \Delta.
\label{eq:extraction_bornee_etat_art}
\end{equation}
C'est l'Hypoth\`ese d'Extraction Born\'ee~\cite{Smith09}, formalis\'ee \`a la
section~\ref{sec:cro_classique} par l'\'equation~\eqref{eq:extraction_bornee} et
fond\'ee par la Proposition~\ref{prop:fondement_extraction}.
\end{sloppypar}

\begin{boiteArbitre}
L'arbitre comprend que $\Delta$ le d\'esigne, puisque ce param\`etre mesure ce qu'il tire
de la preuve au-del\`a du message. Il s'agit peut-\^etre de quelques dizaines de bits, qu'il faudra mesurer.
\end{boiteArbitre}

% ----------------------------------------------------------------------------
\subsection{D\'efinitions pr\'eliminaires du cadre CRO}
\label{subsec:definitions_preliminaires_cro}
% ----------------------------------------------------------------------------
\begin{sloppypar}
Cette section introduit les trois
m\'etriques dans leur formulation issue de la th\'eorie de l'information, formulation qui,
seule, autorise le raisonnement asymptotique du Trilemme CRO.

\subsubsection{Confidentialit\'e, fiabilit\'e, opposabilit\'e}
\label{subsubsec:definitions_preliminaires}

La confidentialit\'e mesure l'indistinguabilit\'e au pire cas~\cite{GoldwasserMicali84} selon la d\'efinition
\begin{equation}
\Conf(\mathcal{P}) = 1 - \max_{m_0,m_1,c} \TV(P_{\Sigma|M=m_0,C=c},
P_{\Sigma|M=m_1,C=c}).
\label{eq:conf_def}
\end{equation}
La fiabilit\'e, qui quantifie la r\'ecup\'eration correcte, vaut
\begin{equation}
\Rel(\mathcal{P}) = 1 - \Pr[\hat{M} \neq M],
\label{eq:rel_def}
\end{equation}
o\`u $\hat{M}$ est l'estimateur optimal. L'opposabilit\'e capture la
r\'eduction relative d'incertitude sur le verdict institutionnel~\cite{MINKA2026CLO}, d\'efinie par
\begin{equation}
\Opp(\mathcal{P}, C) = \frac{I(V;\Sigma|C)}{H(V|C)},
\label{eq:opp_def}
\end{equation}
o\`u $V = \mathsf{Interpret}(\Sigma, C)$ est le verdict produit par
l'interpr\`ete institutionnel.

\begin{remark}[Coexistence de deux notations de l'opposabilit\'e]\label{rem:deux_notations}
Deux notations de l'opposabilit\'e coexistent, renvoyant toutes les deux \`a la m\^eme r\'ealit\'e, mais chacune servant un but pr\'ecis.
La notation de l'\'equation~\eqref{eq:opp_def} sert \`a pr\'eciser le protocole sous revue et le contexte particulier de son examen ou ex\'ecution, c'est une notation purement p\'edagogique, utilis\'ee notamment dans les \'equations~\eqref{eq:opp_helios}, \eqref{eq:opp_zksnark} et \eqref{eq:opp_ecdsa}.
La notation de l'\'equation~\eqref{eq:opp_clo} est celle utilis\'ee dans les d\'emonstrations et autres d\'eveloppements, pour sa sym\'etrie avec celles des deux autres m\'etriques.
\end{remark}

\subsubsection{Relation avec les d\'efinitions cryptographiques standard}
\label{subsubsec:relation_crypto}

Les m\'etriques CRO s'articulent avec les d\'efinitions cryptographiques classiques.

\textbf{S\'emantique s\'ecurit\'e.} La confidentialit\'e $\Conf$ g\'en\'eralise la s\'ecurit\'e s\'emantique de Goldwasser et Micali~\cite{GoldwasserMicali84}, puisque $\Conf = 1$ implique l'indistinguabilit\'e s\'emantique parfaite.

\textbf{Infalsifiabilit\'e existentielle.} La fiabilit\'e de \emph{v\'erification} $\Relver$ (Remarque~\ref{rem:rel_ver}) est reli\'ee \`a la notion d'infalsifiabilit\'e existentielle sous attaque adaptative~\cite{GMR88}, car une signature infalsifiable v\'erifie $\Relver \to 1$. La fiabilit\'e de reconstruction $\Rel$ de l'\'equation~\eqref{eq:rel_def} mesure, elle, la capacit\'e \`a restituer le message.

\textbf{Non-r\'epudiation.} L'opposabilit\'e $\Opp$ n'a pas d'\'equivalent direct dans la litt\'erature cryptographique classique, bien qu'elle soit implicitement vis\'ee par les travaux sur la non-r\'epudiation~\cite{ZhouGollmann}.

\subsubsection{M\'etriques quantiques}
\label{subsubsec:metriques_quantiques}

Les d\'efinitions quantiques g\'en\'eralisent les m\'etriques classiques~\cite{Watrous2009,Boneh2011}.

\textbf{Confidentialit\'e quantique.} On d\'efinit $\Conf^q$ par
\begin{equation}
\Conf^q(\mathcal{P}) = 1 - \frac{1}{2}\max_{m_0,m_1,c}
\|\rho_\Sigma^{m_0,c} - \rho_\Sigma^{m_1,c}\|_1.
\label{eq:conf_q_def}
\end{equation}

\textbf{Fiabilit\'e quantique.} La fiabilit\'e $\Rel^q$ est d\'efinie via l'estimateur POVM optimal
maximisant la probabilit\'e de r\'ecup\'eration correcte~\cite{Wilde2017}.

\textbf{Opposabilit\'e quantique.} L'institution obtient le verdict $V$
\`a partir de l'issue classique d'une mesure sur l'\'etat de preuve et l'on pose
\begin{equation}
\Opp^q(\mathcal{P}) = \frac{I(V;\Sigma|C)}{H(V|C)}.
\label{eq:opp_q_def}
\end{equation}
La borne de Holevo joue ici un r\^ole d\'ecisif, puisque aucune mesure ne permet \`a un interpr\`ete
institutionnel d'extraire de l'\'etat de preuve plus que son information de
Holevo~\cite{Holevo1973}.
\end{sloppypar}

% ----------------------------------------------------------------------------
\subsection{Trilemmes et bornes d'impossibilit\'e}
\label{subsec:trilemmes}
% ----------------------------------------------------------------------------
\begin{sloppypar}
Trois r\'esultats
d'impossibilit\'e ont directement guid\'e la formulation du Trilemme CRO. Tous suivent
le m\^eme sch\'ema, qui consiste \`a identifier une contrainte externe puis \`a d\'emontrer
qu'elle borne irr\'em\'ediablement les compromis r\'ealisables.

\subsubsection{Le th\'eor\`eme d'Arrow}
\label{subsubsec:arrow}

Le th\'eor\`eme d'impossibilit\'e d'Arrow~\cite{Arrow50} \'etablit qu'aucune r\`egle
d'agr\'egation des pr\'ef\'erences individuelles ne peut satisfaire simultan\'ement
cinq conditions rationnelles (universalit\'e, unanimit\'e, indiff\'erence aux
alternatives non pertinentes, souverainet\'e, non-dictature). Ce r\'esultat
fondateur de la th\'eorie du choix social illustre comment des contraintes
apparemment mod\'er\'ees peuvent engendrer des impossibilit\'es profondes~\cite{Sen}.

\subsubsection{Le th\'eor\`eme CAP}
\label{subsubsec:cap}

Le th\'eor\`eme CAP de Brewer~\cite{Brewer00}, prouv\'e formellement par Gilbert
et Lynch~\cite{GilbertLynch02}, \'etablit qu'un syst\`eme distribu\'e ne peut
simultan\'ement garantir coh\'erence (Consistency), disponibilit\'e (Availability)
et tol\'erance aux partitions (Partition tolerance). L'impossibilit\'e est
conditionnelle, puisqu'elle ne s'applique qu'en pr\'esence de partitions r\'eseau~\cite{ChandraHadzilacos96}.

\subsubsection{L'impossibilit\'e FLP}
\label{subsubsec:flp}

Fischer, Lynch et Paterson~\cite{FLP85} ont prouv\'e qu'aucun algorithme
d\'eterministe ne peut r\'esoudre le consensus dans un syst\`eme asynchrone en
pr\'esence d'une seule d\'efaillance. L'impossibilit\'e est conditionn\'ee \`a l'asynchronie~\cite{BenOr83,DebratMerz2012}.

\subsubsection{Le trilemme blockchain}
\label{subsubsec:trilemme_blockchain}

Buterin~\cite{Buterin17} a formalis\'e que les syst\`emes blockchain ne peuvent
simultan\'ement atteindre d\'ecentralisation, s\'ecurit\'e et \'evolutivit\'e
(scalability). Ce trilemme, contrairement \`a CAP et FLP, n'est pas issu d'une preuve
formelle mais d'une heuristique d'ing\'enierie~\cite{Nakamoto08,Garay15,Narayanan16}.

\subsubsection{Impossibilit\'es en cryptographie}
\label{subsubsec:impossibilites_crypto}

Plusieurs r\'esultats d'impossibilit\'e sp\'ecifient les limites fondamentales en cryptographie.

\textbf{Obfuscation.} Barak et al.~\cite{Barak2001} ont prouv\'e l'impossibilit\'e
de l'obfuscation virtuelle black-box pour les programmes g\'en\'eraux.

\textbf{Privacy-Utility trade-off.} La litt\'erature sur la \textit{privacy}
diff\'erentielle~\cite{Dwork} \'etablit des bornes sur les compromis entre
confidentialit\'e et utilit\'e.

\medskip
Ces paradigmes partagent une m\^eme structure, qui consiste \`a identifier une
contrainte externe (partition, asynchronie, d\'ecentralisation) puis \`a
d\'emontrer qu'elle borne les compromis r\'ealisables. Le Trilemme CRO suit ce
sch\'ema avec $\Delta$ comme contrainte dominante~; son extension post-quantique
ajoute $\eta_q$ comme contrainte suppl\'ementaire li\'ee aux effets quantiques.
\end{sloppypar}

% ----------------------------------------------------------------------------
\subsection{V\'erification formelle}
\label{subsec:verification_formelle}
% ----------------------------------------------------------------------------
\begin{sloppypar}
La v\'erification formelle
modifie le statut \'epist\'emique d'une preuve. Une preuve
v\'erifi\'ee m\'ecaniquement garantit l'exactitude de chaque \'etape de calcul et r\'ev\`ele
les hypoth\`eses que la version papier laissait implicites~\cite{Wiedijk2007}.

\subsubsection{Panorama des assistants de preuve}
\label{subsubsec:assistants_preuve}

Coq~\cite{Coq} (Calcul des Constructions Inductives, preuve du Th\'eor\`eme
des Quatre Couleurs~\cite{Gonthier08}, compilateur CompCert~\cite{Leroy09}) et Isabelle/HOL~\cite{Isabelle}
sont matures et largement d\'eploy\'es. Lean4~\cite{Lean4,moura2015lean} combine fondements
logiques rigoureux et programmation fonctionnelle performante, avec la
biblioth\`eque Mathlib~\cite{Mathlib} pour l'appui math\'ematique.

\subsubsection{Formalisation de la cryptographie}
\label{subsubsec:formalisation_crypto}

Plusieurs projets ont formalis\'e des r\'esultats cryptographiques.

\textbf{EasyCrypt.} Barthe et al.~\cite{BartheEA2013} ont d\'evelopp\'e un assistant
de preuve sp\'ecialis\'e pour les preuves cryptographiques computationnelles,
utilisant le calcul de relations probabilistes.

\textbf{F* et HACL*.} Le projet Everest~\cite{BhargavanEA2017,FStar,HACL} vise
une impl\'ementation v\'erifi\'ee de HTTPS, avec la biblioth\`eque cryptographique
HACL* formellement v\'erifi\'ee.

\textbf{CryptoVerif.} Blanchet~\cite{CryptoVerif} a d\'evelopp\'e un outil automatique
pour prouver les propri\'et\'es de s\'ecurit\'e des protocoles cryptographiques.

\textbf{CryptHOL.} Lochbihler~\cite{Lochbihler2016} a formalis\'e les jeux
cryptographiques en Isabelle/HOL.

\textbf{Th\'eorie de l'information avec Coq/Isabelle.} Affeldt et al.~\cite{Affeldt2015}
et Hibon-Portier~\cite{HibonPortier20} ont formalis\'e la th\'eorie de l'information
dans Coq et Isabelle respectivement.

\subsubsection{Lean4 pour le trilemme CRO}
\label{subsubsec:lean4_trilemme}

Lean4 a \'et\'e retenu \`a l'issue d'une comparaison syst\'ematique sur quatre crit\`eres.

\textbf{Couverture math\'ematique.} Mathlib (v4.34.0) fournit l'analyse r\'eelle,
les probabilit\'es et la divergence de Kullback--Leibler
(\texttt{InformationTheory.KullbackLeibler})~\cite{Mathlib}~; les in\'egalit\'es
de Pinsker et de Fano n'y sont pas formalis\'ees et sont donc pos\'ees comme
propri\'et\'es des instances de protocole (section~\ref{sec:lean4}).

\textbf{Proximit\'e syntaxique.} La syntaxe de Lean4 reste lisible par un
math\'ematicien sans formation sp\'ecialis\'ee en v\'erification formelle~\cite{moura2015lean}.

\textbf{Performance de compilation.} Lean4 compile vers C natif et
la v\'erification du noyau alg\'ebrique du Trilemme CRO prend environ une
minute et demie une fois Mathlib t\'el\'echarg\'ee. Le code fourni en
Annexe~\ref{chap:annexe_lean4}, fig\'e sur Lean~v4.34.0 et Mathlib~v4.34.0 dans le
d\'ep\^ot compagnon~\cite{CROTrilemmaDepot}, peut aussi \^etre coll\'e dans
l'\'editeur en ligne \url{https://live.lean-lang.org/}, qui utilise la version
courante de Mathlib.

\textbf{Dynamisme de la communaut\'e.} Croissance forte depuis 2022 avec
contributions actives en cryptographie formelle.

\subsubsection{Apports \'epist\'emiques de la formalisation}
\label{subsubsec:apports_epistemiques}

Au-del\`a de la certification~\cite{AbadiRogaway00,BanaComon12}, la formalisation
Lean4 (section~\ref{sec:lean4}) a produit trois apports.

\begin{enumerate}
\item \textbf{Explicitation d'une hypoth\`ese.} La formalisation a impos\'e de poser
explicitement le r\'egime de pr\'eservation de la confidentialit\'e
$I(M;\Sigma|C) \leq \kappa/\lambda$ (Hypoth\`ese~\ref{ass:regime}), que la preuve
publi\'ee d\'eduisait d'une \'etape d'optimisation. La formulation alternative
retenue le pose comme hypoth\`ese.

\item \textbf{Formulations alternatives.} La formalisation a conduit \`a des
formulations plus pr\'ecises de trois \'enonc\'es publi\'es
(section~\ref{subsec:lean4_erreurs}). La forme raffin\'ee du trilemme, qui
d\'elimite le r\'egime o\`u la capacit\'e institutionnelle est contraignante, a
\'egalement \'et\'e certifi\'ee.

\item \textbf{S\'eparation explicite.} La formalisation s\'epare le noyau alg\'ebrique
v\'erifi\'e des propri\'et\'es pos\'ees (Fano, Pinsker, hypoth\`eses institutionnelles),
clarifiant le statut \'epist\'emique de chaque composant~\cite{Zellagui2018}.
\end{enumerate}
\end{sloppypar}

% ----------------------------------------------------------------------------
%\subsection{Synth\`ese}
%\label{sec:synthese_chap2}
% ----------------------------------------------------------------------------

\begin{boiteArbitre}
L'arbitre referme le dossier. Ce seul chapitre lui a fait traverser six corpus
(th\'eorie de l'information, m\'ecanique quantique, primitives cryptographiques,
sciences cognitives, trilemmes d'impossibilit\'e, v\'erification formelle),
ce qui lui a d'abord paru beaucoup. Il comprend maintenant que chaque pi\`ece \'etait n\'ecessaire.

Ce qui le frappe en premier tient \`a une absence. Aucun des corpus pass\'es en
revue ne nommait l'incapacit\'e qu'il \'eprouve chaque jour \`a extraire
un verdict d'une preuve pourtant irr\'eprochable. La quantit\'e $\Opp = I(V;\Sigma|C)/H(V|C)$
n'apparaissait nulle part. Ce travail l'introduit comme m\'etrique de premi\`ere
classe, ce que l'arbitre re\c{c}oit comme une reconnaissance.

Il note ensuite la coh\'erence de l'\'edifice. Les g\'en\'eralisations quantiques
(Holevo, Pinsker quantique, Fano quantique) s'encha\^inent avec la m\^eme rigueur
que leurs analogues classiques, si bien que le cadre s'\'etend au r\'egime quantique
sans se fracturer.

L'armature analytique repose sur trois in\'egalit\'es dont d\'epend toute la
m\'ecanique du Trilemme CRO. L'in\'egalit\'e de Pinsker relie la confidentialit\'e
\`a l'information mutuelle, celle de Fano borne la fiabilit\'e par la probabilit\'e
d'erreur, tandis que la DPI garantit que le verdict n'invente pas d'information.

Le r\^ole de Lean4 l'\'etonne davantage. Outre la certification, la formalisation
a contraint \`a expliciter une hypoth\`ese et \`a pr\'eciser plusieurs \'enonc\'es. L'arbitre
ne s'attendait pas \`a voir la rigueur formelle servir d'instrument de d\'ecouverte.

Il tourne la page. Six objectifs de recherche l'attendent (Obj~\ref{obj:obj1_unification},~Obj~\ref{obj:obj2_impossibilite},~Obj~\ref{obj:obj3_postquantique},~Obj~\ref{obj:obj4_saturation},~Obj~\ref{obj:obj5_verification} et~Obj~\ref{obj:obj6_experimentation}). Il veut savoir
s'ils tiennent ensemble.
\end{boiteArbitre}